
\SetRowSpace{2mm}
\AnswerColumnBreak
\TypeQuestionDouble{}{construct_vector_operations}
\FigProblemBlock{50mm}{fig/lesson_01_triangle_construction}{\Text{図の三角形 $\mathrm{ABC}$ において、次のベクトルを、始点を $\mathrm{A}$ とする有向線分として図示しなさい。}

\QQ{$\overrightarrow{\mathrm{AB}}+\overrightarrow{\mathrm{AC}}$}
  {\begin{minipage}[t]{0.98\linewidth}%
     \vspace{0pt}%
     \LMFigBlock[\linewidth]{fig/lesson_01_triangle_construction_01_a}%
   \end{minipage}}
\QQ{$\overrightarrow{\mathrm{AB}}-\overrightarrow{\mathrm{AC}}$}
  {\begin{minipage}[t]{0.98\linewidth}%
     \vspace{0pt}%
     \LMFigBlock[\linewidth]{fig/lesson_01_triangle_construction_02_a}%
   \end{minipage}}
\QQ{$2\overrightarrow{\mathrm{AB}}+\overrightarrow{\mathrm{AC}}$}
  {\begin{minipage}[t]{0.98\linewidth}%
     \vspace{0pt}%
     \LMFigBlock[\linewidth]{fig/lesson_01_triangle_construction_03_a}%
   \end{minipage}}
\QQ{$-\overrightarrow{\mathrm{AB}}+2\overrightarrow{\mathrm{AC}}$}
  {\begin{minipage}[t]{0.98\linewidth}%
     \vspace{0pt}%
     \LMFigBlock[\linewidth]{fig/lesson_01_triangle_construction_04_a}%
   \end{minipage}}
}{}



\TypeQuestion{$\vec{a},\vec{b}$ は一次独立なベクトルである。4点 $\mathrm{A},\mathrm{B},\mathrm{C},\mathrm{D}$ の位置ベクトルが $\vec{a},3\vec{b},2\vec{a}+\vec{b},-\vec{a}+2\vec{b}$ であるとき、直線 $\mathrm{AB}$ と直線 $\mathrm{CD}$ の交点を $\mathrm{P}$ とする。}{intersection_from_two_parameter_forms}
\QQ{点 $\mathrm{P}$ が直線 $\mathrm{AB}$ 上にあることを用いて、$\vec{p}$ を実数 $s$ と $\vec{a},\vec{b}$ を用いて表しなさい。}{}
\QQ{点 $\mathrm{P}$ が直線 $\mathrm{CD}$ 上にあることを用いて、$\vec{p}$ を実数 $t$ と $\vec{a},\vec{b}$ を用いて表しなさい。}{}
\QQ{実数 $s,t$ と $\vec{p}$ を求めなさい。}{}
\QQ{点 $\mathrm{P}$ は線分 $\mathrm{AB},\mathrm{CD}$ の内部にあるか、延長上にあるかをそれぞれ答えなさい。}{}



\TypeQuestion{点 $\mathrm{A},\mathrm{B},\mathrm{C}$ の位置ベクトルをそれぞれ $\vec{a},\vec{b},\vec{c}$ とする。面積と点 $\mathrm{P}$ の位置に関する条件から、点 $\mathrm{P}$ の位置ベクトルを表しなさい。}{position_vector_from_area_ratios}
\QQ{$\Tri{ABC}$ の内部に点 $\mathrm{P}$ があり、$\Tri{PBC}:\Tri{PCA}:\Tri{PAB}=3:5:7$ である。}{}
\QQ{$\Tri{ABC}$ の内部に点 $\mathrm{P}$ があり、$\Tri{PAB}$ の面積は $\Tri{ABC}$ の面積の $\dfrac25$、$\Tri{PBC}$ の面積は $\Tri{ABC}$ の面積の $\dfrac14$ である。}{}
\QQ{点 $\mathrm{P}$ は直線 $\mathrm{AC}$ に関して点 $\mathrm{B}$ と反対側にあり、他の2辺の直線については $\Tri{ABC}$ の内部と同じ側にある。また、$\Tri{PBC}:\Tri{PCA}:\Tri{PAB}=2:1:3$ である。}{}
\QQ{点 $\mathrm{P}$ は直線 $\mathrm{BC}$ に関して点 $\mathrm{A}$ と反対側にあり、他の2辺の直線については $\Tri{ABC}$ の内部と同じ側にある。また、$\Tri{PBC}:\Tri{PCA}:\Tri{PAB}=1:4:2$ である。}{}


